\chapter{Python Environment Setup}

This appendix provides instructions for setting up the Python environment to run all code examples in this book.

\section{System Requirements}

\begin{itemize}
    \item Python 3.8 or higher
    \item pip (Python package installer)
    \item Git (for cloning the repository)
\end{itemize}

\section{Installation}

\subsection{Using pip}

\begin{lstlisting}[style=bash, caption={Installing required packages}]
pip install numpy scipy matplotlib pytest
\end{lstlisting}

\subsection{Using conda}

\begin{lstlisting}[style=bash, caption={Installing with conda}]
conda create -n gauss-book python=3.10
conda activate gauss-book
conda install numpy scipy matplotlib pytest
\end{lstlisting}

\subsection{Development Installation}

\begin{lstlisting}[style=bash, caption={Installing in development mode}]
git clone https://github.com/yourusername/gauss-book.git
cd gauss-book
pip install -e .
\end{lstlisting}

\section{Running the Code}

\subsection{Running Individual Modules}

Each chapter has a corresponding Python module in the \texttt{python/} directory:

\begin{lstlisting}[style=bash]
# Chapter 1: Modular Arithmetic
python -m python.modular

# Chapter 2: Quadratic Reciprocity
python -m python.quadratic

# Chapter 3: Gauss Sums
python -m python.gauss_sums

# Chapter 7: AGM
python -m python.agm

# Chapter 8: Gaussian Elimination
python -m python.gaussian_elim

# Chapter 9: Least Squares
python -m python.least_squares

# Chapter 10: Gauss Quadrature
python -m python.gauss_quadrature

# Chapter 11: Normal Distribution
python -m python.normal

# Chapter 12: Gaussian Processes
python -m python.gp

# Chapter 13: Heptadecagon
python -m python.heptadecagon

# Chapter 14: Theorema Egregium
python -m python.theorema

# Chapter 15: Orbital Mechanics
python -m python.orbital

# Additional modules
python -m python.lemniscate
python -m python.gauss_composition
\end{lstlisting}

\subsection{Running Tests}

\begin{lstlisting}[style=bash]
# Run all tests
pytest tests/

# Run specific test file
pytest tests/test_modular.py -v

# Run with coverage
pytest --cov=python tests/
\end{lstlisting}

\subsection{Using Jupyter Notebooks}

For interactive exploration, we recommend Jupyter:

\begin{lstlisting}[style=bash]
pip install jupyter
jupyter notebook
\end{lstlisting}

Then navigate to the \texttt{notebooks/} directory (if available) or create your own notebooks importing the modules:

\begin{lstlisting}[style=python]
# In a notebook cell
import sys
sys.path.append('../python')

from modular import chinese_remainder, primitive_root
from quadratic import legendre_symbol, tonelli_shanks
from agm import pi_brent_salamin, elliptic_K

# Run examples
print(chinese_remainder([2, 3, 2], [3, 5, 7]))
print(primitive_root(17))
print(pi_brent_salamin(3))
\end{lstlisting}

\section{Dependencies Summary}

\begin{table}[H]
\centering
\begin{tabular}{ll}
\toprule
\textbf{Package} & \textbf{Version} \\
\midrule
numpy & $\ge$ 1.21 \\
scipy & $\ge$ 1.7 \\
matplotlib & $\ge$ 3.4 \\
pytest & $\ge$ 6.2 \\
Pillow & $\ge$ 8.0 (for image generation) \\
\bottomrule
\end{tabular}
\end{table}

\section{Troubleshooting}

\begin{itemize}
    \item \textbf{ModuleNotFoundError}: Ensure you're running from the project root directory or have added the \texttt{python/} directory to your PYTHONPATH.
    \item \textbf{Numerical precision issues}: Some algorithms (AGM, Gaussian quadrature) converge very quickly; adjust tolerance parameters if needed.
    \item \textbf{Import errors}: Make sure all dependencies are installed: \texttt{pip install -r requirements.txt}
\end{itemize}

\section{requirements.txt}

\begin{lstlisting}[style=txt, caption={requirements.txt}]
numpy>=1.21.0
scipy>=1.7.0
matplotlib>=3.4.0
pytest>=6.2.0
Pillow>=8.0.0
\end{lstlisting}